\section*{第 \ref{ch:g11:func} 章 --- 函数与变化}

\begin{solution}{exo:g11:func:1}
$f$：需 $3 - x \geq 0$，定义域为 $\intoc{-\infty}{3}$。

$g$：需 $x^2 - 4 \neq 0$，即 $x \neq \pm 2$，定义域为
$\R \setminus \{-2, 2\}$。

$h$：需 $x \geq 0$ 且 $x \neq 5$，定义域为
$\intco{0}{5} \cup \intoo{5}{+\infty}$。
\end{solution}

\begin{solution}{exo:g11:func:2}
$\sqrt{11} < \sqrt{13}$，因为 $\sqrt{}$ 递增，且 $11 < 13$。

$(-2.1)^2 > (-2.05)^2$，因为 $x^2$ 在
$\intoc{-\infty}{0}$ 上递减，且 $-2.1 < -2.05$。

$\frac{1}{0.9} > \frac{1}{1.1}$，因为 $\frac1x$ 在
$\intoo{0}{+\infty}$ 上递减，且 $0.9 < 1.1$。

$(-1.01)^3 < (-0.99)^3$，因为 $x^3$ 递增，且 $-1.01 < -0.99$。
\end{solution}

\begin{solution}{exo:g11:func:3}
$f(-x) = 3(-x)^4 - (-x)^2 = 3x^4 - x^2 = f(x)$，为偶函数。

$g(-x) = -x^3 - x = -g(x)$，为奇函数。

$h(1) = 2$，且 $h(-1) = 1 - 1 = 0$，故 $h(-1)$ 既不等于 $h(1)$，
也不等于 $-h(1)$，因此既非偶函数，也非奇函数。

$k(-x) = \frac{-x}{x^2+1} = -k(x)$，为奇函数。
\end{solution}

\begin{solution}{exo:g11:func:4}
由变化表可知，在 $\intcc{-3}{1}$ 上，$f$ 由 $-2$ 增至 $4$，
且恰好一次取到值 $1$；在 $\intcc{1}{5}$ 上，由 $4$ 降至 $0$，
又恰好一次取到 $1$：方程 $f(x) = 1$ 有
$2$ 个解。$f$ 的最大值为 $f(1) = 4 < 5$，故方程 $f(x) = 5$ 无
解。
\end{solution}

\begin{solution}{exo:g11:func:5}
设 $2 \leq u < v$，则
\[
f(v) - f(u) = v^2 - 4v - u^2 + 4u = (v-u)(v+u) - 4(v-u)
= (v - u)(u + v - 4).
\]
两个因式均为正：$v - u > 0$，且 $u + v > 2 + 2 = 4$，因为
$u \geq 2$，$v > 2$。于是 $f(v) > f(u)$，故 $f$ 在
$\intco{2}{+\infty}$ 上严格递增。
\end{solution}

\begin{solution}{exo:g11:func:6}
$f = \sqrt{u}$，其中 $u(x) = 5 - x$ 在
$\intoc{-\infty}{5}$ 上递减且非负，故 $f$ 在此递减
（见\cref{prop:g11:func:transforms}）。

$g = \frac1u$，其中 $u(x) = x^2 + 1$ 在
$\intco{0}{+\infty}$ 上递增且为正，故 $g$ 在此递减。

$h = 3 - 2\sqrt x$：$\sqrt x$ 递增，乘以 $-2 < 0$
使单调性反转，加上 $3$ 不改变单调性：$h$ 在
$\intco{0}{+\infty}$ 上递减。
\end{solution}

\begin{solution}{exo:g11:func:7}
\emph{1.} $f$ 为顶点式：在 $\intoc{-\infty}{2}$ 上递减，
在 $\intco{2}{+\infty}$ 上递增，最小值为 $f(2) = -3$。

\emph{2.} $g = \frac{1}{u}$，其中 $u(x) = (x-2)^2 + 1 = f(x) + 4$。
平移 $4$ 不改变 $f$ 的单调性，且 $u \geq 1 > 0$。
取倒数则使其反转：$g$ 在 $\intoc{-\infty}{2}$ 上递增，且
在 $\intco{2}{+\infty}$ 上递减，最大值为
$g(2) = \frac{1}{1} = 1$。
\end{solution}

\begin{solution}{exo:g11:func:8}
任何经过这三点的山丘形曲线均可。则：
$f + 2$ 经过 $(0,3)$、$(2,5)$、$(4,3)$，单调性不变
（竖直平移）；$-f$ 经过 $(0,-1)$、$(2,-3)$、$(4,-1)$，
单调性\emph{相反}（山峰变为谷底）；$2f$ 经过
$(0,2)$、$(2,6)$、$(4,2)$，单调性不变，但沿竖直方向
放大。
\end{solution}

\begin{solution}{exo:g11:func:9}
$2 - \frac{1}{x+1} = \frac{2(x+1) - 1}{x+1} = \frac{2x+1}{x+1} = f(x)$。
在 $\intoo{-1}{+\infty}$ 上，$u(x) = x + 1$ 递增且为正，故
$\frac1u$ 递减，$-\frac1u$ 递增，加上 $2$ 不改变
单调性，故 $f$ 在 $\intoo{-1}{+\infty}$ 上严格递增。
\end{solution}

\begin{solution}{exo:g11:func:10}
\emph{1. 真命题。} 若 $u < v$，则
$(f+g)(v) - (f+g)(u) = \bigl(f(v)-f(u)\bigr) + \bigl(g(v)-g(u)\bigr)$ 为
两个非负数之和。

\emph{2. 假命题。} $f(x) = g(x) = x$ 在 $\R$ 上递增，但它们的
积 $x^2$ 在 $\R$ 上并不递增（它在
$\intoc{-\infty}{0}$ 上递减）。（若两个函数还都
非负，则该命题成立。）

\emph{3. 假命题。} 同一反例：$f(x) = x$ 在 $\R$ 上递增，
但 $f^2(x) = x^2$ 并不递增。
\end{solution}

\begin{solution}{exo:g11:func:11}
\emph{1.} 对 $0 < u < v$，有
\[
f(v) - f(u) = (v - u) + \frac1v - \frac1u
= (v - u) + \frac{u - v}{uv}
= (v - u)\left(1 - \frac{1}{uv}\right)
= (v - u)\,\frac{uv - 1}{uv}.
\]

\emph{2.} 恒有 $v - u > 0$，$uv > 0$。若 $0 < u < v \leq 1$，则
$uv < 1$，故 $uv - 1 < 0$，且 $f(v) < f(u)$，因此 $f$ 在
$\intoc{0}{1}$ 上严格递减。若 $1 \leq u < v$，则 $uv > 1$，且 $f(v) > f(u)$：
$f$ 在 $\intco{1}{+\infty}$ 上严格递增。

\emph{3.} $f$ 在 $1$ 之前递减，之后递增，故
$f(x) \geq f(1) = 2$ 对一切 $x > 0$ 成立，等号当且仅当 $x = 1$ 时取得。
（这一不等式以多种面貌反复出现；参看
\cref{ex:g11:quad:canonical} 以及
\cref{ch:g12:deriv} 中的切线不等式。）
\end{solution}

\begin{solution}{pb:g11:func:1}
\textbf{1.} $x^4 - 3x^2$：偶函数（$(-x)^4 - 3(-x)^2 = x^4 -
3x^2$）。$x^3 - x$：奇函数。$x^2 + x$：非奇非偶（$f(-1) = 0 \neq
f(1) = 2$ 且 $\neq -f(1)$）。$\abs x$：偶函数。$\frac1x$：奇函数。

\textbf{2.} 偶函数：图像关于 $y$ 轴对称
（如镜中成像）。奇函数：图像关于原点对称（半圈旋转）。

\textbf{3.} 若 $f$、$g$ 均为奇函数：
$(fg)(-x) = f(-x)g(-x) = (-f(x))(-g(x)) = (fg)(x)$，故为偶函数。
列表：偶 $\times$ 偶 $=$ 偶；偶 $\times$ 奇 $=$ 奇；
奇 $\times$ 奇 $=$ 偶 --- 即奇偶性上的符号法则。

\textbf{4.} 奇函数在 $0$ 处：$f(0) = -f(0)$，故 $f(0) = 0$。偶函数：
$f(-3) = f(3) = 7$。既是偶函数又是奇函数：
$f(x) = f(-x) = -f(x)$ 对一切 $x$ 成立，故 $2f(x) = 0$，$f$ 为
零函数。

\textbf{5.} 右半部分（$x \geq 0$）完全决定一个
偶函数：左半部分即其镜像。既非偶函数也非奇函数：
$x^2 + x$（第 1 题）——其两种对称性在
$x = 1$ 处均不成立。

\textbf{6.} $E(x) = \frac{(x^3 + x^2 + 1) + (-x^3 + x^2 +
1)}{2} = x^2 + 1$（偶函数），$O(x) = \frac{(x^3 + x^2 + 1) -
(-x^3 + x^2 + 1)}{2} = x^3$（奇函数）；两者之和即为 $f$。

\textbf{7.} $E(-x) = \frac{f(-x) + f(x)}{2} = E(x)$，故为偶函数。
$O(-x) = \frac{f(-x) - f(x)}{2} = -O(x)$，故为奇函数。又
由直接相加得 $E(x) + O(x) = f(x)$：每个函数
都可如此分解。

\textbf{8.} 由 $E_1 + O_1 = E_2 + O_2$ 得
$E_1 - E_2 = O_2 - O_1$。左边为偶函数，右边为
奇函数 --- 故该函数既是偶函数又是奇函数，从而为零函数（第 4 题）：
$E_1 = E_2$，且 $O_1 = O_2$。分解唯一。

\textbf{9.} $E(x) = \frac12\left(\frac{1}{1 - x} + \frac{1}{1 +
x}\right) = \frac12 \cdot \frac{(1 + x) + (1 - x)}{1 - x^2}
= \frac{1}{1 - x^2}$，且
$O(x) = \frac12 \cdot \frac{(1 + x) - (1 - x)}{1 - x^2}
= \frac{x}{1 - x^2}$。验证：$E$ 为偶函数，$O$ 为奇函数，且
两者之和为 $\frac{1 + x}{1 - x^2} =
\frac{1 + x}{(1 - x)(1 + x)} = \frac{1}{1 - x} = f(x)$。

\textbf{10.} 由于每一部分都服从某种对称性，
每一部分只需一半的工作（及一半的数据）：在一侧求解，
再镜像到另一侧 --- 而许多运算（积分、
级数、信号分析）处理对称的部分远比处理原始函数
简便。

\textbf{11.} $(x-3)^2$：向右平移 $3$ 个单位。$x^2 + 2$：向上平移
$2$ 个单位。$(x-3)^2 + 2$：两者兼有。$-x^2$：上下翻转。
$2x^2$：沿纵向拉伸 $2$ 倍（开口更窄）。

\textbf{12.} 当 $x \geq 2$ 时：$(x - 2) + (x + 2) = 2x$。当
$-2 \leq x \leq 2$ 时：$(2 - x) + (x + 2) = 4$。当 $x \leq -2$ 时：
$-2x$。图像以斜率 $-2$ 下降，
随后在高度 $4$ 处形成平底，位于 $-2$ 与 $2$ 之间，又以斜率 $2$ 上升：
一个带平底的山谷形。

\textbf{13.} 竖直平移：$f(x) + 2$ 仍为偶函数
（$f(-x) + 2 = f(x) + 2$）。竖直伸缩：$2f(x)$ 仍为
偶函数。水平平移：$(x - 3)^2$ 不是偶函数
（$f(-1) = 16 \neq 4 = f(1)$）：横向滑动破坏了
镜像对称。

\textbf{14.} 在平底上取值为 $4 \neq 6$。当
$x \geq 2$ 时：$2x = 6$，$x = 3$。当 $x \leq -2$ 时：$-2x = 6$，
$x = -3$。解为 $\pm 3$。

\textbf{15.} 设 $f(x) = a(x - 2)^2 - 3$，且 $f(0) = 4a - 3 = 1$，
则 $a = 1$，故 $f(x) = (x - 2)^2 - 3$。

\textbf{16.} $g(v) - g(u) = (v - u) + 4\,\frac{u - v}{uv}
= (v - u)\,\frac{uv - 4}{uv}$：当 $u < v$ 且
$uv < 4$ 时为负，$uv > 4$ 时为正：于是在
$\intoc{0}{2}$ 上递减，在 $\intco{2}{+\infty}$ 上递增；最小值在
$x = 2$ 处取得，值为 $g(2) = 4$。

\textbf{17.} 两边长分别为 $x$ 与 $\frac{16}{x}$：
$P(x) = 2\left(x + \frac{16}{x}\right)$，取最小值（第
16 题的方法，将 $4 \to 16$，转折点 $\sqrt{16} = 4$）
在 $x = 4$ 处取得，即 $4 \times 4$ 的\emph{正方形}，周长为 $16$ ---
狄多的答案，现已对一切实数边长证明成立。

\textbf{18.} 在 $\intco{0}{+\infty}$ 上：$x \mapsto x^2 + 1$
递增，$\sqrt{\phantom{x}}$ 也递增，故
复合函数递增。又 $h(-x) = \sqrt{x^2 + 1} = h(x)$：$h$
为偶函数，由镜面对称可知它在
$\intoc{-\infty}{0}$ 上递减。最小值为 $h(0) = 1$。

\textbf{19.} 点 $(3, 0)$ 到点
$(x, \sqrt x)$ 的距离平方为
$d^2(x) = (x - 3)^2 + x = x^2 - 5x + 9$，在
$x = \frac52$ 处取得最小值。最近点
$\left(\frac52, \sqrt{\frac52}\right)$，距离
$\sqrt{\frac{11}{4}} = \frac{\sqrt{11}}{2}$。求 $d^2$ 的最小值
是合理的，因为平方在
正数上严格递增：$d$ 与 $d^2$ 在同一处取得最小值。

\textbf{20.} 原料：须熟记于心的基本曲线。
加工机器：平移、翻转与伸缩，用以由旧图像组装新图像；
质量控制：奇偶性与奇偶分解，
它们使每项研究减半。测量台：因式分解
$f(v) - f(u)$，由此直接读出变化与最小值。下一章的
有力工具是\emph{导数}，它能
一眼读出这些性质。
\end{solution}
